% from main.tex
%%%%%%%%%%  scan idx 146  |  folio 133  %%%%%%%%%%
Soient \(X\in\Ob S\), et \(m,n\in\mathbb Z\) (\(m\leq n\)). Pour que \(F\in\Ob\Dpx^-({}_A X)\) soit d'amplitude plate contenue dans \([m,n]\) il faut et il suffit qu'il existe une famille couvrante \((X_i\to X)\) telle que \(F|X_i\) soit d'amplitude plate contenue dans \([m,n]\). En particulier, si \(S/X\) est de type fini (SGA 4 VI 1.7), et si \(F\) est localement de tor-dimension finie, alors \(F\) est de tor-dimension finie.

\textbf{Proposition 5.6.} \emph{Sous les hypothèses de (2.16), si \(E'\) est d'amplitude plate contenue dans \([m',n']\) en tant qu'objet de \(\Dpx^-({}_A, S')\) et si \(E\) est d'amplitude plate contenue dans \([m,n]\), alors \(E'\overset{\mathrm L}{\otimes}_A E\) est d'amplitude plate contenue dans \([m+m',n+n']\).}

\emph{Preuve.} C'est immédiat : utiliser la forme (5.1 (i)) de la notion de tor-dimension finie et le fait que, si \(E\) est un \(A\)-Module à gauche plat, \(f^{-1}(E)\) est un \(f^{-1}(A)\)-Module à gauche plat (SGA 4 IV App. II).

En particulier :

\textbf{Corollaire 5.6.1.} \textbf{a)} Le foncteur
\[
  \mathrm L f^*:\Dpx^-({}_A S)\longrightarrow \Dpx^-({}_A, S')
\]
induit un foncteur
\[
  \mathrm L f^*:\Dpx({}_A S)_{\torf}\longrightarrow \Dpx({}_A, S')_{\torf}.
\]

\textbf{b)} Si \(A\) est commutatif, le foncteur
\[
  \overset{\mathrm L}{\otimes}_A:
  \Dpx^-(S)\times_S\Dpx^-(S)
  \longrightarrow \Dpx^-(S)
\]
induit un foncteur
\[
  \overset{\mathrm L}{\otimes}_A:
  \Dpx(S)_{\torf}\times_S\Dpx(S)_{\torf}
  \longrightarrow \Dpx(S)_{\torf}.
\]

\textbf{Remarque 5.6.2.} Soient \(E\in\Ob\Dpx^-({}_A S)\) (\(S\)= objet final de \(S\)), et \(m,n\in\mathbb Z\) (\(m\leq n\)). Si \(E\) est d'amplitude plate contenue dans \([m,n]\), il en est de même, d'après (5.6), de \(E_s\) pour tout point \(s\) de \(S\). La réciproque est vraie si \(S\) a assez de points (utiliser la forme (5.1.(ii)) de la notion d'amplitude
